\section*{Capítulo \ref{ch:b3:solid-state} --- Química del estado sólido: bandas, defectos y semiconductores}

\begin{solution}{exo:b3:solid-state:1}
$(E - \alpha)/|\beta| = -2\cos(k\pi/7)$: $-1.802$, $-1.247$, $-0.445$, $+0.445$, $+1.247$, $+1.802$. Abarcan
$3.604|\beta|$, ya el 90\,\% de la anchura de banda $4|\beta|$ de la cadena infinita.
\end{solution}

\begin{solution}{exo:b3:solid-state:2}
$2N$ ($N$ orbitales, con dos electrones cada uno). El sodio, con un electrón 3s por átomo, llena
a medias su banda s: es un metal. El magnesio tiene dos y llenaría exactamente la banda s,
pero las bandas 3s y 3p se solapan, de modo que los niveles ocupados más altos pertenecen a una
banda combinada parcialmente llena: también es un metal.
\end{solution}

\begin{solution}{exo:b3:solid-state:3}
$\lambda \approx \qty{1239.8}{eV.nm}/E_g$: fosfuro de galio, \qty{549}{nm} (verde); nitruro de galio, \qty{365}{nm}
(ultravioleta cercano). Los diodos azules usan nitruro de galio aleado con indio, cuya
\omterm{def:b3:solid-state:band}{banda prohibida} es más estrecha.
\end{solution}

\begin{solution}{exo:b3:solid-state:4}
$\ce{CaCl2} \xrightarrow{\ce{NaCl}} \mathrm{Ca_{Na}^{\bullet} + V_{Na}' + 2\,Cl_{Cl}^{\times}}$: dos posiciones catiónicas y dos aniónicas, como en el NaCl; un Ca y
dos Cl; carga $+1 - 1 = 0$. $\ce{Y2O3} \xrightarrow{\ce{ZrO2}} \mathrm{2\,Y_{Zr}' + 3\,O_O^{\times} + V_O^{\bullet\bullet}}$: dos posiciones catiónicas y cuatro de oxígeno; dos Y
y tres O; carga $-2 + 2 = 0$.
\end{solution}

\begin{solution}{exo:b3:solid-state:5}
$T^{3/2} = 5196$ a \qty{300}{K}, $kT = \qty{0.02585}{eV}$. Silicio: $\sqrt{N_cN_v} = \qty{2.42e19}{cm^{-3}}$, $\eu^{-1.125/0.05170} = \num{3.55e-10}$,
$n_i = \qty{8.6e9}{cm^{-3}}$. Germanio: $\sqrt{N_cN_v} = \qty{7.16e18}{cm^{-3}}$, $\eu^{-0.661/0.05170} = \num{2.80e-6}$, $n_i = \qty{2.0e13}{cm^{-3}}$. Cociente:
$\num{4.3e-4}$: la \omterm{def:b3:solid-state:band}{banda prohibida} más estrecha del germanio le da más de dos mil veces más
portadores.
\end{solution}

\begin{solution}{exo:b3:solid-state:6}
$n = N_D = \qty{1.0e16}{cm^{-3}}$; $p = n_i^2/n = \num{1.0e20}/\num{1.0e16} = \qty{1.0e4}{cm^{-3}}$.
\end{solution}

\begin{solution}{exo:b3:solid-state:7}
$E_F - E_i = kT\ln(n/n_i) = \qty{0.02585}{eV} \times \ln(10^6) = \qty{0.36}{eV}$.
\end{solution}

\begin{solution}{exo:b3:solid-state:8}
$kT = \qty{0.06894}{eV}$ a \qty{800}{K}; $n/N = \eu^{-2.3/0.1379} = \num{5.7e-8}$; en un mol, $\num{6.02e23} \times \num{5.7e-8} = \num{3.4e16}$
pares.
\end{solution}

\begin{solution}{exo:b3:solid-state:9}
La masa de una celda es $\rho a^3 = \qty{5.70}{g.cm^{-3}} \times (\qty{4.300e-8}{cm})^3 = \qty{4.532e-22}{g}$, es decir,
$\qty{4.532e-22}{g} \times N_A = \qty{272.9}{g}$ por mol de celdas, o \qty{68.23}{g} por oxígeno. Entonces
$(1 - x) \times 55.8 + 16.0 = 68.23$ da $1 - x = 0.936$: $x = 0.064$, $\mathrm{Fe_{0.936}O}$.
\end{solution}

\begin{solution}{exo:b3:solid-state:10}
$\ln(\sigma T)$: $-3.51$, $-1.30$, $0.36$ para $1/T = 3.333$, $2.857$, $\qty{2.500e-3}{K^{-1}}$; pendiente entre los puntos extremos:
$(0.36 + 3.51)/(-\qty{8.33e-4}{K^{-1}}) = -\qty{4.65e3}{K}$ (el punto central está sobre la recta), luego
$E_a = \qty{4.65e3}{K} \times \qty{8.617e-5}{eV.K^{-1}} = \qty{0.40}{eV}$.
\end{solution}

\begin{solution}{exo:b3:solid-state:11}
$\tfrac12\ce{O2} \rightleftharpoons \mathrm{O_O^{\times} + V_M' + h^{\bullet}}$, $K = [\mathrm{V_M'}][\mathrm h^\bullet]/p(\ce{O2})^{1/2}$. Balance de carga: $[\mathrm h^\bullet] = [\mathrm{V_M'}]$, luego
$[\mathrm h^\bullet]^2 = K\,p(\ce{O2})^{1/2}$ y $[\mathrm h^\bullet] \propto p(\ce{O2})^{1/4}$; la conductividad sigue a los \omterm{def:b3:solid-state:carriers}{huecos}.
\end{solution}

\begin{solution}{exo:b3:solid-state:12}
$n/N = \sqrt{N_i/N}\,\eu^{-\Delta H_F/2kT} = \sqrt 2\,\eu^{-1.4/(2 \times 0.05170)} = 1.414 \times \num{1.32e-6} = \num{1.9e-6}$.
\end{solution}

\begin{solution}{pb:b3:solid-state:1}
\textbf{1.} $\ce{Y2O3} \xrightarrow{\ce{ZrO2}} \mathrm{2\,Y_{Zr}' + 3\,O_O^{\times} + V_O^{\bullet\bullet}}$.

\textbf{2.} Posiciones: dos posiciones catiónicas y cuatro de oxígeno (tres llenas y una vacía),
la razón 1\,:\,2 del $\ce{ZrO2}$. Masa: 2 Y y 3 O a cada lado. Carga: $2 \times (-1) + 2 = 0$.

\textbf{3.} Una vacante por cada dos iones itrio: media por itrio.

\textbf{4.} Por cada $(\ce{ZrO2})_{0.92}(\ce{Y2O3})_{0.08}$ hay 0.92 Zr y 0.16 Y, 1.08 cationes: $x = 0.16/1.08 = 0.148$,
$\mathrm{Zr_{0.852}Y_{0.148}O_{1.926}}$.

\textbf{5.} $(x/2)/2 = 0.0370$: el 3.7\,\% de las posiciones de oxígeno están vacías.

\textbf{6.} $8 \times 0.0370 = 0.296$ vacantes por celda.

\textbf{7.} $a^3 = (\qty{5.14e-8}{cm})^3 = \qty{1.358e-22}{cm^3}$: $0.296/\qty{1.358e-22}{cm^3} = \qty{2.2e21}{cm^{-3}}$.

\textbf{8.} Cada ion óxido debe saltar por encima de una barrera de \qty{1.0}{eV} hacia una vacante
vecina; la fracción de intentos que tienen éxito, $\eu^{-E_a/kT}$, crece muy deprisa con $T$.

\textbf{9.} A \qty{1073.15}{K}, $kT = \qty{0.09248}{eV}$: $A = \sigma T\eu^{E_a/kT} = 0.030 \times 1073.15 \times \eu^{10.81} = \qty{1.6e6}{S.K.cm^{-1}}$.

\textbf{10.} A \qty{973.15}{K}: $\sigma = (A/T)\eu^{-11.92} = \qty{0.011}{S.cm^{-1}}$.

\textbf{11.} A \qty{573.15}{K}: $\sigma = (A/T)\eu^{-20.25} = \qty{4.5e-6}{S.cm^{-1}}$.

\textbf{12.} $R = L/(\sigma S) = \qty{0.10}{cm}/(\sigma \times \qty{0.20}{cm^2})$: \qty{46}{\ohm} a \qty{700}{\celsius}, $\qty{1.1e5}{\ohm}$ a \qty{300}{\celsius}.

\textbf{13.} $R < \qty{1}{k\ohm}$ exige $\sigma > \qty{5.0e-4}{S.cm^{-1}}$; resolver $(A/T)\eu^{-E_a/kT} = \num{5.0e-4}$ (por tanteo)
da $T \approx \qty{760}{K}$, unos \qty{490}{\celsius}.

\textbf{14.} El escape está frío tras el arranque y a baja carga; el calefactor lleva el
electrolito a su temperatura de funcionamiento en pocos segundos, de modo que el control
funciona desde el arranque, cuando se emite la mayor parte de los contaminantes.

\textbf{15.} $\ce{O2 + 4e- -> 2O^2-}$; en \omterm{def:b3:solid-state:kroger-vink}{notación de Kröger--Vink}, $\ce{O2} + \mathrm{2\,V_O^{\bullet\bullet} + 4\,e'} \longrightarrow \mathrm{2\,O_O^{\times}}$.

\textbf{16.} El oxígeno se reduce en el lado del aire, y los iones óxido se oxidan de nuevo
a oxígeno en el lado del escape: la celda transfiere $\ce{O2}$ del aire al escape, con
$\Delta_rG = RT\ln(p_{\mathrm{esc}}/p_{\mathrm{aire}})$ por mol de $\ce{O2}$ y cuatro electrones: $E = -\Delta_rG/4F = (RT/4F)\ln(p_{\mathrm{aire}}/p_{\mathrm{esc}})$.

\textbf{17.} $RT/4F = 8.3145 \times 973.15/(4 \times 96485) = \qty{0.02096}{V}$.

\textbf{18.} Cero: las dos presiones son iguales.

\textbf{19.} $(RT/4F)\ln 10 = \qty{48.3}{mV}$ por década.

\textbf{20.} $\qty{0.02096}{V} \times \ln(0.2095/0.010) = \qty{0.02096}{V} \times 3.04 = \qty{0.064}{V}$.

\textbf{21.} $\qty{0.02096}{V} \times \ln(0.2095/\num{1.0e-20}) = \qty{0.02096}{V} \times 44.49 = \qty{0.93}{V}$.

\textbf{22.} El escape rico contiene monóxido de carbono, hidrógeno e hidrocarburos sin
quemar; sobre el electrodo de platino, el poco oxígeno que queda reacciona con ellos,
y $p(\ce{O2})$ la fijan los equilibrios $\ce{2CO + O2 <=> 2CO2}$ y $\ce{2H2 + O2 <=> 2H2O}$, que dejan un valor
ínfimo a \qty{700}{\celsius}.

\textbf{23.} $p = 0.2095\,\eu^{-0.45/0.02096} = \qty{1.0e-10}{bar}$.

\textbf{24.} En torno a la razón estequiométrica, el escape pasa de un ligero exceso
de oxígeno a un ligero exceso de CO e hidrógeno: $p(\ce{O2})$ cae unas dieciocho
potencias de diez para un pequeño cambio de combustible, y la tensión salta de menos de
\qty{0.1}{V} a unos \qty{0.9}{V}. El control del motor añade combustible cuando la tensión es baja y
lo retira cuando es alta, manteniendo la mezcla en el punto estequiométrico, en el que el
catalizador de tres vías convierte a la vez el CO, los hidrocarburos y los óxidos de nitrógeno.

\textbf{25.} Con un escape rico a \qty{700}{\celsius}, la sonda da unos \qty{0.93}{V}.
\end{solution}
